\section{Nonlinear departure on the exact invariant level}
\label{sec:r12-nonlinear}

We now turn the real spectral pair into an actual nonlinear trajectory.
The equation throughout this section is the pure mass-critical quartic
model
\begin{equation}\label{eq:r12-pure-quartic}
 iu_t=\Delta^2u-|u|^{8/N}u,\qquad N\in\{1,2,3\},
 \qquad X=H^2(\mathbb R^N;\mathbb C).
\end{equation}
The space $X$ is regarded as real. In this concrete phase space, the
branch, weighted regularity, and complementary gap used in the spectral theorem in Appendix~\ref{sec:r11-spectral} follow from real-profile nondegeneracy. That remaining
assumption is stated explicitly below. The form decomposition is the
one-negative-direction instance of the Hamiltonian trichotomy developed
by Lin and Zeng~\cite{LZ}; the estimates needed here are proved directly.

\begin{theorem}[Ancient departure and nonlinear orbital instability]
\label{thm:r12-nonlinear-instability}
Let $Q$ be a real sharp ground profile at frequency one. Assume
\begin{equation}\label{eq:r12-real-nondegeneracy}
 \ker\left(\Delta^2+1-\left(1+\frac8N\right)|Q|^{8/N}\right)
   =\operatorname{span}_{\mathbb R}
        \{\partial_1Q,\ldots,\partial_NQ\}.
\end{equation}
Then there is a local $C^2$ profile family $U_{\omega,v}$ through $Q$,
unique up to phase and translation near $Q$. Put $U_v=U_{1,v}$.
For every sufficiently small fixed $v\ne0$, put
\[
 \mathcal O_v=\{e^{i\theta}U_v(\cdot-y):\theta\in\mathbb R,
                                     \ y\in\mathbb R^N\}.
\]
There is a nonstationary solution in the co-moving, co-rotating frame,
defined for $t\le0$, of the form
\begin{equation}\label{eq:r12-ancient-asymptotic}
 w_v(t)=U_v+\varepsilon_v e^{\lambda_vt}Y_{+,v}
              +O_v(e^{(1+\beta)\lambda_vt}),\qquad t\to-\infty,
\end{equation}
where $\|Y_{+,v}\|_{H^2}=1$, $J\mathcal L_vY_{+,v}
=\lambda_vY_{+,v}$, $\varepsilon_v>0$, and
\[
 \beta=\begin{cases}1,&N=1,2,\\2/3,&N=3.\end{cases}
\]
At every time this solution has exactly the same mass, momentum, and
energy as $U_v$.

Consequently, there are $\rho_v>0$ and data $u_{\delta,v}$ satisfying
\begin{gather*}
 M(u_{\delta,v})=M(U_v),\quad P(u_{\delta,v})=P(U_v),\quad
 E(u_{\delta,v})=E(U_v),\\
 u_{\delta,v}=U_v+\delta Y_{+,v}+O_v(\delta^{1+\beta})
                         \quad\hbox{in }H^2,
\end{gather*}
whose first exit time from the orbital distance $\rho_v$ obeys
\begin{equation}\label{eq:r12-exit-time}
 \tau_{\delta,v}=\lambda_v^{-1}\log(1/\delta)+O_v(1)
 \qquad(\delta\downarrow0).
\end{equation}
The spectral calculation gives
\begin{equation}\label{eq:r12-rate}
 \lambda_v=\sqrt{\frac{2v^TTv}{\mathcal D_Q}}+o(|v|),
\end{equation}
where
\begin{gather*}
 T_{ij}=\langle\partial_iQ,L_-^{-1}\partial_jQ\rangle,
 \qquad c_j=\langle\Lambda Q,L_-^{-1}\partial_jQ\rangle,\\
 \mathcal D_Q=\langle\Lambda Q,L_-^{-1}\Lambda Q\rangle-c^TT^{-1}c.
\end{gather*}
In particular, $e^{it}U_v(x-vt)$ is orbitally unstable in $H^2$ even
on its exact joint mass--momentum--energy level.
\end{theorem}

\subsection{Local flow and the nonlinear remainder}

Set $F(z)=|z|^{8/N}z$, with $\mathbb C$ identified with $\mathbb R^2$.
Since $H^2\hookrightarrow L^\infty\cap W^{1,4}$ for $N\le3$ and
$H^2$ is an algebra, $F:X\to X$ is locally $C^{1,\beta}$ with the
exponent in the theorem. Here is a direct check of the regularity that
is needed. Twice differentiating $DF(w)h$ produces sums of
\begin{gather*}
 DF(w)\nabla^2h,\qquad D^2F(w)[\nabla w,\nabla h],\\
 D^2F(w)[\nabla^2w,h],\qquad D^3F(w)[\nabla w,\nabla w,h].
\end{gather*}
The products lie in $L^2$ by the stated embeddings. On bounded subsets
of $\mathbb R^2$, $D^3F$ is Lipschitz in dimensions one and two, and
$2/3$-H\"older in dimension three. Subtracting these expressions at
$w_1,w_2$ therefore gives
\[
 \|DF(w_1)-DF(w_2)\|_{\mathcal B(X)}
       \le C_R\|w_1-w_2\|_X^\beta
\]
for $\|w_i\|_X\le R$ and $\|w_1-w_2\|_X\le1$.
The lower derivative terms obey the same bound. Smooth approximation
justifies the weak chain rules. It follows that, near a fixed $U$,
\begin{align}
 \|F(U+h)-F(U)-DF(U)h\|_X&\le C\|h\|_X^{1+\beta},
 \label{eq:r12-remainder}\\
 \|R(h_1)-R(h_2)\|_X
 &\le C(\|h_1\|_X^\beta+\|h_2\|_X^\beta)\|h_1-h_2\|_X,
 \label{eq:r12-remainder-lipschitz}
\end{align}
where $R(h)=i[F(U+h)-F(U)-DF(U)h]$.

The free group $e^{-it\Delta^2}$ is unitary on $X$. A contraction
argument gives local well-posedness and continuation while the $H^2$
norm is bounded. Mass, momentum, and energy are conserved. For a direct
justification, insert a selfadjoint Fourier cutoff $P_n$ into the
Hamiltonian nonlinearity $P_nF(P_nu)$, differentiate the conserved
quantities for this regularized flow, and pass to the $H^2$ solution
uniformly on compact existence intervals. The cutoff commutes with
phase and translation actions. Thus no Strichartz or spectral-mapping
hypothesis is required for the local nonlinear flow.

For
\[
 w(t,x)=e^{-it}u(t,x+vt)=U_v+h(t,x)
\]
the perturbation equation is
\begin{equation}\label{eq:r12-perturbation-eq}
 h_t=\mathcal A_vh+R_v(h),\qquad
 \mathcal A_v=J\mathcal L_v,\quad J=-i.
\end{equation}
The phase and translation transformations are $H^2$ isometries and
preserve all three conserved quantities.

\subsection{Why the other spectral assumptions follow here}

\begin{lemma}[Local branch and complementary gap]
\label{lem:r12-branch}
For a real sharp quartic ground profile satisfying
\eqref{eq:r12-real-nondegeneracy}, all the other assumptions of the
small-velocity theorem in Appendix~\ref{sec:r11-spectral} hold. In particular,
$\Lambda Q\in H^4$, the local profile branch is $C^2$ into $H^4$,
and its only Hessian zero modes are phase and translations.
\end{lemma}
\begin{proof}
Sharp phase coercivity gives $L_-\ge0$, $\ker L_-=\mathbb RQ$,
and an inverse on $Q^\perp$. The nonlinearity calculation above and
$(\Delta^2+1)Q=F(Q)$ give $Q\in H^6$. Therefore
$V=|Q|^{8/N}\in W^{2,\infty}$ and $V(x)\to0$ as $|x|\to\infty$.
We first justify the scaling derivative without assuming a moment.

Choose smooth bounded truncations $\phi_R$ of
$a\langle x\rangle$, with $\phi_R\to a\langle x\rangle$ pointwise,
$\phi_R\le Ca\langle x\rangle$, and
$|\nabla\phi_R|+|\Delta\phi_R|\le Ca$ uniformly in $R$.
Test $(\Delta^2+1-V)Q=0$ with $e^{2\phi_R}Q$ and put
$w_R=e^{\phi_R}Q$. The exact real identity
\[
 \operatorname{Re}\langle\Delta^2Q,e^{2\phi_R}Q\rangle
 =\|\Delta w_R+|\nabla\phi_R|^2w_R\|_2^2
  -\|2\nabla\phi_R\cdot\nabla w_R+(\Delta\phi_R)w_R\|_2^2
\]
shows that the difference from $\|\Delta w_R\|_2^2$ is bounded by
$Ca^2\|w_R\|_{H^2}^2$ for $a\le1$.
Choose a fixed ball outside which $V\le1/4$. Inside the ball,
$\int V|w_R|^2$ is uniformly bounded because the weights there are
bounded independently of $R$. Taking $a>0$ small absorbs both the
exterior potential and the $Ca^2$ error. Thus $w_R$ is uniformly
bounded in $H^2$, and passage to the limit gives
$w=e^{a\langle x\rangle}Q\in H^2$.

The conjugated equation
\[
 e^{a\langle x\rangle}(\Delta^2+1)e^{-a\langle x\rangle}w=Vw
\]
has smooth bounded coefficients in its lower-order terms, of order at
most three. Since $V\in W^{2,\infty}$, elementary elliptic bootstrapping
improves $w$ successively from $H^2$ to $H^3,H^4,H^5,H^6$.
At the first step, the third-order terms belong to $H^{-1}$; at the
subsequent steps they belong to $L^2,H^1,H^2$, respectively.
Consequently $\Lambda Q=NQ/2+x\cdot\nabla Q\in H^4$, and
$L_+\Lambda Q=-4Q$ follows by differentiating the scaled equation.

For the branch, set
\[
 \mathcal F(U,\omega,v)=\Delta^2U+\omega U+iv\cdot\nabla U-F(U).
\]
This map is $C^2$ from $H^4\times\mathbb R^{1+N}$ to $L^2$.
Let $g_0=iQ$, $g_j=\partial_jQ$, and augment the profile equation by
\[
 \mathcal F(U,\omega,v)=\sum_{j=0}^N b_jg_j,
 \qquad \langle U-Q,g_j\rangle_{\mathbb R}=0.
\]
Its derivative with respect to $(U,b)$ is invertible: the real
selfadjoint Hessian at $Q$ is Fredholm, its kernel is exactly the
span of the $g_j$, and it has a bounded complementary inverse from
$L^2$ to $H^4$. The implicit function theorem gives a $C^2$ augmented
branch. Phase and translation invariance imply
\[
 \langle\mathcal F(U,\omega,v),iU\rangle_{\mathbb R}=0,
 \qquad
 \langle\mathcal F(U,\omega,v),\partial_kU\rangle_{\mathbb R}=0.
\]
The pairing matrix between the fixed $g_j$ and these moving symmetry
vectors is close to the invertible Gram matrix at $Q$. Hence $b=0$,
and the branch consists of exact profiles.

The gauge determines profiles uniquely near $Q$ up to phase and
translation. Mass-preserving scaling maps parameters to
$(\kappa^4\omega,\kappa^3v)$; after returning to the gauge, uniqueness
shows that the branch is scaling-closed. Its reduced action is $C^2$
and has the usual charge derivatives. Differentiation at $(1,0)$ gives
$D_{v_j}U=-iL_-^{-1}\partial_jQ$ in the fixed gauge.

Finally, the Hessians have common domain $H^4$. The potentials converge
in $L^\infty$, and $iv\cdot\nabla$ is small from $H^4$ to $L^2$.
Norm-resolvent convergence therefore preserves the rank $N+1$ of a
small isolated spectral cluster about zero. The $N+1$ independent
vectors $iU,\partial_jU$ fill that cluster. They belong to the operator
domain by the profile equation and elliptic bootstrapping. Thus there
are no further zero or near-zero eigenvalues, and the inverse on the
kernel complement, extended by zero, is norm-continuous on $L^2$.
This verifies the complementary-gap assumption rather than adding it.
\end{proof}

\subsection{The linear group bound from one negative direction}

\begin{lemma}\label{lem:r12-group-bound}
For each fixed sufficiently small $v\ne0$,
\[
 n_-(\mathcal L_v)=1,
 \qquad
 \|e^{t\mathcal A_v}\|_{\mathcal B(H^2)}\le C_v e^{\lambda_vt}
 \quad(t\ge0).
\]
On an invariant complement to the two real eigenvectors the group
has at most linear growth.
\end{lemma}
\begin{proof}
At $v=0$, sharp critical optimality gives nonnegativity of $L_+$ on
$Q^\perp$ and
\[
 \langle Q,L_+Q\rangle=-\frac8N\|Q\|_{2+8/N}^{2+8/N}<0.
\]
Since $L_-\ge0$, the full Hessian has exactly one negative direction.
For small $v$, the free operator
\[
 H_v=\Delta^2+1+iv\cdot\nabla
\]
has symbol $|\xi|^4+1-v\cdot\xi$ uniformly comparable to
$1+|\xi|^4$. The potential $DF(U_v)$ is a decaying bounded real
symmetric matrix and converges in $L^\infty$ to $DF(Q)$. The drift
is small from $H^4$ to $L^2$. This gives norm-resolvent convergence
with a common lower bound and a uniformly positive essential spectrum.
The isolated zero cluster contains exactly the prescribed $N+1$
symmetry modes. Hence $n_-(\mathcal L_v)=1$ for small $v$.
This concrete perturbation argument excludes negative spectrum arriving
from infinity, as well as from zero.

Fix $v$ and omit its subscript. The profile equation and $F(U)\in H^2$
give $U\in H^6$. The matrix multiplier $DF(U)$ is bounded on $H^2$.
The free operator $-iH_v$ generates an isometric group on $H^2$;
its bounded real-linear perturbation $iDF(U)$ therefore generates a
$C^0$ group there. The $L^2$ eigenvectors of $\mathcal A$ lie in
$H^4$, and their eigenvalue equations improve them to $H^6$.

Let $\mathcal A Y_\pm=\pm\lambda Y_\pm$, $\lambda>0$, and write
$B(f,g)=\langle\mathcal Lf,g\rangle$ for the real form. The identity
\[
 B(\mathcal Af,g)+B(f,\mathcal Ag)=0
\]
gives $B(Y_\pm,Y_\pm)=0$. Furthermore
$b=B(Y_+,Y_-)\ne0$: otherwise their span would be a two-dimensional
totally isotropic space modulo $\ker\mathcal L$, impossible for a
form of negative index one. Their span is disjoint from the kernel,
as is seen by applying $\mathcal A$ and $\mathcal A^2$ to any
proposed relation. Hence
\[
 W=\operatorname{span}_{\mathbb R}\{Y_+,Y_-\},\qquad
 Z=W^{\perp_B},\qquad X=W\oplus Z
\]
is an invariant decomposition with bounded projection
\[
 \Pi_Wh=\frac{B(h,Y_-)}bY_++\frac{B(h,Y_+)}bY_-.
\]
The form is nonnegative on $Z$, because its sole negative direction
is already in $W$. Its null space in $Z$ is exactly
$K=\ker\mathcal L$: if $B(z,z)=0$, nonnegativity implies
$B(z,Z)=0$, while the definition gives $B(z,W)=0$.

Let $\Pi_K$ be the $L^2$ orthogonal projection onto $K$. Compactness
of the decaying potential relative to the positive free form gives
\begin{equation}\label{eq:r12-center-coercivity}
 B(z,z)\ge c_v\|(I-\Pi_K)z\|_{H^2}^2,
 \qquad z\in Z.
\end{equation}
Indeed, a violating sequence in $Z\cap K^\perp$, normalized in $H^2$,
has a weak limit. Weak lower semicontinuity of the free form and
compactness of the potential force its form value to be nonpositive.
The preceding null-space statement makes that limit zero. Compactness
then makes the potential term vanish, contradicting coercivity of the
free form and the normalization.

The form $B(e^{t\mathcal A}z,e^{t\mathcal A}z)$ is conserved, first
on the generator domain and then by density. It controls the component
orthogonal to $K$ through \eqref{eq:r12-center-coercivity}. For an
$L^2$ orthonormal basis $g_j$ of $K$, write
$h=z+k$ with $k=\Pi_Kh$. In the $H^{-2}$--$H^2$ pairing,
\[
 \frac d{dt}\langle h,g_j\rangle
 =-\langle\mathcal Lz,Jg_j\rangle,
 \qquad
 \left|\frac d{dt}\langle h,g_j\rangle\right|
       \le C_v\|z\|_{H^2}.
\]
This proves
\[
 \|e^{t\mathcal A}|_Z\|_{\mathcal B(H^2)}\le C_v(1+|t|).
\]
The group on $W$ is exactly $e^{\pm\lambda t}$. Absorbing $1+t$
into a $\lambda$-dependent multiple of $e^{\lambda t}$ proves the
lemma.
\end{proof}

\subsection{An ancient nonlinear orbit and exact conservation}

\begin{proof}[Proof of Theorem~\ref{thm:r12-nonlinear-instability}]
Lemma~\ref{lem:r12-branch} verifies the hypotheses of the calculation in Appendix~\ref{sec:r11-spectral} and hence supplies $\lambda_v$ and
$Y_{+,v}$. Fix $v$ and suppress it. Normalize $\|Y_+\|_X=1$. On continuous
functions $h:(-\infty,0]\to X$ use the norm
\[
 \|h\|_* =\sup_{t\le0}e^{-\lambda t}\|h(t)\|_X.
\]
For small $\varepsilon>0$, consider on the ball $\|h\|_*\le2\varepsilon$
the map
\begin{equation}\label{eq:r12-ancient-fixedpoint}
 \Phi(h)(t)=\varepsilon e^{\lambda t}Y_+
       +\int_{-\infty}^t e^{(t-s)\mathcal A}R(h(s))\,ds.
\end{equation}
The integral converges in $X$, since Lemma~\ref{lem:r12-group-bound}
and \eqref{eq:r12-remainder} give
\begin{equation}\label{eq:r12-ancient-error}
 \|\Phi(h)(t)-\varepsilon e^{\lambda t}Y_+\|_X
 \le \frac{C_v C_R2^{1+\beta}}{\beta\lambda}
             \varepsilon^{1+\beta}e^{(1+\beta)\lambda t}.
\end{equation}
By \eqref{eq:r12-remainder-lipschitz}, its weighted Lipschitz constant
is at most $C'_v\varepsilon^\beta/(\beta\lambda)$.
Choose $\varepsilon$ small enough for invariance of the ball and a
contraction constant below $1/2$. The fixed point is a nonlinear mild
solution on every finite subinterval: split the integral in
\eqref{eq:r12-ancient-fixedpoint} at a finite initial time and use
the group law. It satisfies \eqref{eq:r12-ancient-asymptotic} and
converges strongly to $U$ backward in time. Conservation and continuity
of $M,P,E$ therefore imply
\begin{equation}\label{eq:r12-exact-invariants}
 (M,P,E)(U+h(t))=(M,P,E)(U),\qquad t\le0.
\end{equation}

It remains to establish departure from the entire symmetry orbit.
The map $(\theta,y)\mapsto e^{i\theta}U(\cdot-y)$ is $C^2$ into
$H^2$, with injective differential and tangent $K$. It has no
nontrivial stabilizer: a nonzero phase-periodic translation contradicts
$U\in L^2$. Nor can translations tending to infinity approach $U$ in
$L^2$. Thus every orbit point sufficiently near $U$ belongs to this
local chart. Put
\[
 \eta_v=\operatorname{dist}_{H^2}(Y_+,K)>0.
\]
It is positive because $\mathcal A$ annihilates $K$ and has nonzero
eigenvalue on $Y_+$. Taylor expansion of the chart gives, for small
$a>0$,
\begin{equation}\label{eq:r12-orbit-transverse}
 \operatorname{dist}_{H^2}(U+aY_+,\mathcal O_v)
                          \ge\frac{\eta_v}{2}a.
\end{equation}
Indeed, any competing orbit point with comparable distance lies in the
chart with parameters $O(a)$; projecting its expansion off $K$ proves
the inequality. No existence of a global nearest orbit point is needed.

Shrink $\varepsilon_v$ further so that
\eqref{eq:r12-orbit-transverse} holds at $a=\varepsilon_v$, and the
right side of \eqref{eq:r12-ancient-error} at $t=0$ is at most
$\eta_v\varepsilon_v/4$. Then
\[
 \operatorname{dist}_{H^2}(U+h(0),\mathcal O_v)
                         \ge\eta_v\varepsilon_v/4.
\]
Set $\rho_v=\eta_v\varepsilon_v/8$. For $0<\delta<\varepsilon_v$
let $t_\delta=\lambda^{-1}\log(\delta/\varepsilon_v)$ and
$u_{\delta,v}=U+h(t_\delta)$. Its expansion and exact invariant values
follow from \eqref{eq:r12-ancient-error} and
\eqref{eq:r12-exact-invariants}. Its solution in the moving frame is
$U+h(t+t_\delta)$. The first time $\tau_{\delta,v}$ at which the
orbital distance equals $\rho_v$ satisfies
\[
 \lambda^{-1}\log\frac{\rho_v}{2\delta}
 \le\tau_{\delta,v}\le
 \lambda^{-1}\log\frac{\varepsilon_v}{\delta},
\]
because $\|h(t+t_\delta)\|_X\le2\delta e^{\lambda t}$ and the
distance at $t=-t_\delta$ exceeds $\rho_v$.
This proves \eqref{eq:r12-exit-time}. Formula \eqref{eq:r12-rate}
is the asymptotic in Theorem~\ref{thm:r11-spectral-pair} with $m=4$.
\end{proof}

\begin{remark}[Velocity dependence and the remaining assumption]
The construction gives a fixed positive escape radius for each fixed
nonzero small velocity. It does not bound that radius uniformly as
$v\to0$. The eigenvectors approach the neutral sector and the spectral
projection constants may diverge. The rigorously implied rate statement
along a unit direction $e$ is the iterated limit
\[
 \lim_{r\downarrow0}r\left[
   \lim_{\delta\downarrow0}
      \frac{\tau_{\delta,re}}{\log(1/\delta)}\right]
       =\sqrt{\frac{\mathcal D_Q}{2e^TTe}}.
\]
Critical scaling transfers the theorem to every frequency $\omega>0$,
with normalized velocity $v\omega^{-3/4}$ and the exact rate covariance
$\lambda(\omega,v)=\omega\lambda(1,v\omega^{-3/4})$.
The nonlinear instability mechanism is a classical Hamiltonian one;
the present connection uses the previously computed critical
small-velocity eigenvalue. Real-profile nondegeneracy is still an
explicit hypothesis. The theorem does not assert finite-time blowup or
classify the subsequent evolution after departure.
\end{remark}
